% =====================================================================
%  GROUP REPORT TEMPLATE -- Applied Groundwater Modelling
% =====================================================================
%  Read REPORT_BRIEF.md first: it says what each section is FOR and how
%  it is judged. This file is only the skeleton.
%
%  Same skeleton in three formats -- use whichever you like:
%      report_template.md    Markdown (JupyterHub, VS Code, any editor)
%      report_template.tex   this file (LaTeX / Overleaf)
%      report_template.docx  Microsoft Word
%
%  HOW TO USE
%    1. Copy this file into your group folder as report.tex.
%    2. Delete every % comment block as you go -- they are instructions,
%       not content.
%    3. Compile to report.pdf in your group folder root.
%       Overleaf: upload and hit Recompile.
%       On the course JupyterHub (pdflatex/xelatex/lualatex are installed):
%           pdflatex report.tex && pdflatex report.tex
%       (run it twice so the table/figure references resolve)
%
%  LIMITS
%    10 pages max, excluding title, references and appendix.
%    Appendix capped at 5 pages.
%    The per-section budgets below add to about 8, which leaves you room.
%
%  Deliberately plain: article class and standard packages only, so it
%  compiles anywhere without installing anything.
% =====================================================================

\documentclass[11pt,a4paper]{article}

\usepackage[utf8]{inputenc}
\usepackage[T1]{fontenc}
\usepackage[margin=2.5cm]{geometry}
\usepackage{graphicx}
\usepackage{booktabs}
\usepackage{hyperref}

% >>> REPLACE EVERY <...> BELOW. A forgotten title block is visible on page 1. <<<
% Two-member group: delete the third \and entry.
\title{<Your title --- name the concession and the question,
        not ``Groundwater Modelling Report''>}
\author{Group <N> \\
        <First Last> \and <Second Last> \and <Third Last>}
\date{<date>}

\begin{document}
\maketitle

% ---------------------------------------------------------------------
\section{Problem and modelling question}
% ~0.5 page. One paragraph.
% Your concession, its setting, and the question your scenario actually
% asks. Be specific: "does the assigned recharge reduction change river
% leakage enough to matter for the concession's yield" is a question;
% "we study the Limmat aquifer" is not.

<...>

% ---------------------------------------------------------------------
\section{Conceptual model and assumptions}
% ~1.5 pages.
% The aquifer as YOU are modelling it: boundaries and what they represent,
% the stresses, your scenario's physical story, and the transport setup --
% spill, doublet, the locked physics (alpha_L = 10 m, alpha_T = 1 m,
% n_e = 0.20) and your group's pinned reactions.
%
% State the assumptions YOU make, not the ones the software makes, and name
% the CONSEQUENCE of each.

<...>

\subsection{Assumptions and their consequences}

\begin{table}[h]
\centering
\begin{tabular}{lll}
\toprule
Assumption & Why reasonable here & What it costs us \\
\midrule
<e.g. steady state> & <...> & <...> \\
<...>               & <...> & <...> \\
\bottomrule
\end{tabular}
\end{table}

% ---------------------------------------------------------------------
\section{Method}
% ~1 page.
% Enough that a competent reader could repeat it: the four flow states and
% what each isolates, how your scenario forcing is applied, the transport
% build, and which quantities came from which export. Point to the
% notebooks; do not reproduce their code.

<...>

% ---------------------------------------------------------------------
\section{Results}
% ~3 pages. THE STRUCTURE MATTERS: organise by FINDING, not by card.
% The full record of each card and its extension is already in the scratch
% notebooks and in tables/ -- do not re-narrate it. What belongs here is the
% INTEGRATED result: what the flow response, the scenario forcing, the water
% balance and the transport verdict say TOGETHER about your concession.
%
% Every figure and table: caption, units, and the file under figures/ or
% tables/ it came from. Every computed number traceable to the frozen bundle.
%
% Two traps to handle explicitly:
%  - head change has two conventions: abs() (counts the injection mound) vs
%    drawdown (base - compare, positive where head falls). Say which.
%  - an advective travel time is not a concentration peak. Do not reconcile
%    them by arithmetic.
%
% Both things are required AT ONCE, and this is the hardest instruction in the brief:
% the section must be finding-shaped, AND every member's work must be visible in it.
% A finding usually draws on more than one card -- say whose analysis established it,
% in the text or the caption, rather than giving each member a subsection.
%
% Rename these headings to YOUR findings.

\subsection{What the doublet does to the flow field}
<...>

\subsection{What your scenario forcing adds, and where it dominates}
<...>

\subsection{What the water balance says}
<...>

\subsection{The transport verdict}
<...>

% Figure pattern -- caption names the source file:
% \begin{figure}[h]
%   \centering
%   \includegraphics[width=0.8\textwidth]{figures/cardA_drawdown_map.png}
%   \caption{Drawdown (base $\rightarrow$ wells), m. Source:
%            \texttt{figures/cardA\_drawdown\_map.png}.}
% \end{figure}

% ---------------------------------------------------------------------
\section{Uncertainty and the limits of your claims}
% ~1.5 pages. This section separates reports at this level. Required:
%  - the mesh envelope on your transport peak (two defensible meshes move it
%    by roughly +/-20 %), and its consequence -- WHETHER OR NOT it flips your
%    verdict;
%  - a signal / artefact / instability judgment on your main flow result with
%    the evidence that discriminates -- or a correct finding that your
%    evidence CANNOT discriminate, and what would;
%  - at least one limitation you found and quantified yourselves;
%  - your card extensions' defensibility answers, synthesised not listed.

<...>

\subsection{What this model cannot tell you}
% State these explicitly. The model resolves transport ALONG the flow axis
% but not across it, so:
%   CLAIMABLE     -- arrival time, first threshold exceedance, receptor peak
%   NOT CLAIMABLE -- contaminated-area maps, lateral threshold-contour width,
%                    exact plume footprint (numerical artefacts at any
%                    feasible grid).
% Then check nothing elsewhere in your report quietly breaks this.

<...>

% ---------------------------------------------------------------------
\section{Conclusion and practical recommendation}
% ~0.5 page.
% What the result means for the actual problem. ONE recommendation, and the
% condition under which you would withdraw it.
%
% A no-reach or below-threshold result is not a weaker result.

<...>

% ---------------------------------------------------------------------
\section{Contribution statement}
% Must be CONSISTENT WITH SUBMISSION_README.md and the STUDENT_NAME fields in the
% scratch notebooks. That file lists one row per card, this one per member --
% different shapes, same people. In a two-member group one member has three cards;
% mark the short one (short).

\begin{table}[h]
\centering
\begin{tabular}{lll}
\toprule
Member & Cards & Contribution \\
\midrule
<...> & <...> & <...> \\
<...> & <...> & <...> \\
<...> & <...> & <...> \\
\bottomrule
\end{tabular}
\end{table}

\noindent\textbf{Steward roles:} flow <name> $\cdot$ transport <name> $\cdot$ export <name>

% ---------------------------------------------------------------------
\section{References}
<...>

% ---------------------------------------------------------------------
\appendix
% Not counted in the 10 pages. Capped at 5 pages.

\section{Figure and table index}

\begin{table}[h]
\centering
\begin{tabular}{ll}
\toprule
Report item & File \\
\midrule
Figure 1 & \texttt{figures/<...>} \\
Table 1  & \texttt{tables/<...>} \\
\bottomrule
\end{tabular}
\end{table}

\section{Derivations or intermediate numbers you want on record}
<...>

\end{document}
